\documentclass{article}
\usepackage{amsmath,amssymb}
\usepackage{booktabs}
\begin{document}

\section{Fit-Uni Algorithm: Principles of Adaptive Weighting and Group Constraint}

\subsection{Core Idea}

Fit-Uni (UniLasso) is a two-stage univariate-guided sparse regression algorithm. The core idea is:
\begin{enumerate}
\item \textbf{Stage 1}: Fit univariate regression for each feature individually, compute leave-one-out (LOO) fitted values, and transform the original features into a new space.
\item \textbf{Stage 2}: Solve the Lasso path with regularization in the transformed feature space.
\end{enumerate}

This article elaborates on the working principles of the two key innovations --- \textbf{adaptive weighting} and \textbf{group constraint}.

\subsection{Adaptive Weighting}

\subsubsection{Step 1: Compute Significance Weights}

Based on univariate regression results, compute significance weights for each feature. Let there be $p$ features, the algorithm:

\begin{enumerate}
\item Extract significance statistic $s_j$:
\begin{itemize}
\item \texttt{correlation}: $s_j = |\mathrm{corr}(X_j, y)|$
\item \texttt{t\_statistic}: $s_j = |t_j|$ ($t_j$ is the t-statistic of the coefficient)
\item \texttt{p\_value}: $s_j = -\log_{10}(p_j)$ ($p_j$ is the p-value of the coefficient)
\end{itemize}

\item Normalize to the interval $[1, w_{\text{max}}]$:
\[
\boxed{
w_j = 1 + \frac{s_j - \min(s)}{\max(s) - \min(s)} \cdot (w_{\text{max}} - 1)
}
\]
The default value is $w_{\text{max}} = 10$.
\end{enumerate}

\subsubsection{Step 2: Role in Proximal Gradient Descent}

In the proximal operation, adaptive weighting directly scales the soft threshold.

For coefficient $w$ after gradient descent update, the thresholds are:
\[
\begin{aligned}
\tau^+ &= \text{lr} \cdot \lambda \cdot w_j \\
\tau^- &= \text{lr} \cdot (\lambda + \lambda_{\text{neg}}) \cdot w_j
\end{aligned}
\]
where:
\begin{itemize}
\item $\text{lr}$: learning rate
\item $\lambda$: L1 regularization strength
\item $\lambda_{\text{neg}}$: extra penalty for negative coefficients (negative penalty)
\end{itemize}

Soft thresholding operation:
\[
\text{prox}(w) =
\begin{cases}
w - \tau^+ & \text{if } w > \tau^+ \\
w + \tau^- & \text{if } w < -\tau^- \\
0 & \text{otherwise}
\end{cases}
\]

\subsubsection{Effect}

\begin{itemize}
\item \textbf{High significance features}: large $w_j$ $\Rightarrow$ large thresholds $\tau^+, \tau^-$ $\Rightarrow$ coefficients are harder to be compressed to zero $\Rightarrow$ more likely to be retained in the model
\item \textbf{Low significance features}: small $w_j$ $\Rightarrow$ small thresholds $\tau^+, \tau^-$ $\Rightarrow$ coefficients are easier to be compressed to zero $\Rightarrow$ more likely to be eliminated
\end{itemize}

Adaptive weighting implements the principle that \textbf{higher significance implies smaller penalty}, improving the consistency of variable selection.

\subsection{Group Constraint}

Group constraint addresses the feature selection problem with collinear features: when multiple features are highly correlated, it encourages them to be selected consistently and maintains sign consistency.

\subsubsection{Step 1: Greedy Correlation Grouping}

\begin{enumerate}
\item Compute the feature correlation matrix $C \in \mathbb{R}^{p \times p}$, where $C_{ij} = \mathrm{corr}(X_i, X_j)$.

\item Greedy grouping algorithm:

\vspace{0.5em}
\noindent\textbf{Algorithm: Greedy Correlation Grouping}
\begin{description}
\item[Input:] Correlation matrix $C \in \mathbb{R}^{p \times p}$, correlation threshold $\rho$, maximum group size $\text{max\_size}$
\item[Output:] List of groups, each group contains feature indices
\end{description}

\begin{enumerate}
\item[1.] Initialize: $\text{assigned} \leftarrow \mathbf{0}_p$, $\text{groups} \leftarrow []$
\item[2.] While not all features are assigned:
    \begin{enumerate}
    \item[$a.$] $importance_j \leftarrow \sum_{i=1}^p |C_{ij}|$
    \item[$b.$] $\text{leader} \leftarrow \arg\max_{j: \text{assigned}_j = 0} importance_j$
    \item[$c.$] $\text{correlated} \leftarrow \{j: \text{assigned}_j = 0 \text{ and } |C_{\text{leader}, j}| \geq \rho\}$
    \item[$d.$] If $|\text{correlated}| > \text{max\_size}$: sort correlated by $|C_{\text{leader}, j}|$ descending, keep top $\text{max\_size}$
    \item[$e.$] Add correlated to groups
    \item[$f.$] Mark correlated as assigned
    \end{enumerate}
\item[3.] Return groups
\end{enumerate}

\vspace{0.5em}
where $\rho = \text{corr\_threshold}$ (default 0.7) is the grouping threshold.
\end{enumerate}

\subsubsection{Step 2: Compute Dominant Sign and Confidence}

For each group $G$:

\begin{enumerate}
\item Weighted voting determines the dominant sign:
\[
\text{vote} = \sum_{j \in G} \mathrm{sign}(\beta_j) \cdot |\beta_j| \cdot w_j
\]
where $\beta_j$ is the univariate regression coefficient, and $w_j$ is the adaptive weight.

\item Dominant sign:
\[
\boxed{
s_G = \mathrm{sign}(\text{vote})
}
\]

\item Sign consistency confidence:
\[
\boxed{
c_G = \frac{|\text{vote}|}{\sum_{j \in G} |\mathrm{sign}(\beta_j) \cdot |\beta_j| \cdot w_j|} \in [0, 1]
}
\]

\item Group penalty weight:
\[
\boxed{
g_j = 1 + 4 \cdot c_G \quad \forall j \in G
}
\]
The range is $[1, 5]$, higher confidence gives larger group penalty.
\end{enumerate}

\subsubsection{Step 3: Asymmetric Proximal Penalty in Gradient Descent}

Combining adaptive weighting and group constraint, the final thresholds are:

\[
\boxed{
\begin{aligned}
\tau^+ &= \text{lr} \cdot \lambda \cdot w_j +
\begin{cases}
0 & \text{if } s_G = +1 \\
\text{lr} \cdot \lambda_G \cdot g_j & \text{if } s_G = -1
\end{cases} \\[1em]
\tau^- &= \text{lr} \cdot (\lambda + \lambda_{\text{neg}}) \cdot w_j +
\begin{cases}
\text{lr} \cdot \lambda_G \cdot g_j & \text{if } s_G = +1 \\
0 & \text{if } s_G = -1
\end{cases}
\end{aligned}
}
\]

where:
\begin{itemize}
\item $\lambda_G = \text{group\_penalty}$: global group penalty strength (default 5.0)
\item $s_G = group\_signs[j]$: group dominant sign
\item $g_j = group\_weights[j]$: group weight
\end{itemize}

\subsubsection{Effect}

\begin{itemize}
\item If group dominant sign is $+1$:
  \begin{itemize}
  \item Positive coefficients: threshold unchanged, easier to retain
  \item Negative coefficients: additional penalty, easier to be compressed to zero
  \end{itemize}
\item If group dominant sign is $-1$:
  \begin{itemize}
  \item Negative coefficients: threshold unchanged, easier to retain
  \item Positive coefficients: additional penalty, easier to be compressed to zero
  \end{itemize}
\end{itemize}

Through asymmetric thresholding, group constraint encourages all coefficients within a group to be consistent with the dominant sign, and inhibits coefficients with inconsistent signs. This effectively reduces false discoveries and promotes grouped selection among highly correlated features.

\subsection{Complete Algorithm Flow Summary}

\begin{center}
\begin{tabular}{lcc}
\hline
Stage & Operation & Output \\
\hline
1 & Input validation and zero-variance filtering & $X_{\text{clean}}, y$ \\
2 & Univariate LOO fitting & $\hat{Z}$ (loo\_fits), $\beta$ (beta\_coefs\_fit) \\
3 & Configure regularization path & $\lambda_1 > \lambda_2 > \dots > \lambda_m$ \\
4 & Compute adaptive weights & $w_j \in [1, w_{\text{max}}]$ for $j=1\dots p$ \\
5 & Correlation grouping & groups, $s_G$ (group\_signs), $g_j$ (group\_weights) \\
6 & Proximal gradient descent solution & $\theta_j(\lambda_k)$ for $k=1\dots m$ \\
7 & Recover coefficients to original space & $\gamma_j(\lambda_k) = \theta_j(\lambda_k) \cdot \beta_j$ \\
\hline
\end{tabular}
\end{center}

Coefficient recovery formula in original space:
\[
\boxed{
\gamma_j(\lambda) = \theta_j(\lambda) \cdot \beta_j
}
\]
Intercept recovery:
\[
\boxed{
\gamma_0(\lambda) = \theta_0(\lambda) + \sum_{j=1}^p \theta_j(\lambda) \cdot \beta_{0j}
}
\]
where $\beta_{0j}$ is the intercept term from univariate regression.

\subsection{Design Motivation}

\begin{center}
\begin{tabular}{ll}
\hline
Problem & Solution \\
\hline
Lasso tends to select unimportant features & Adaptive weighting: higher significance $\Rightarrow$ smaller penalty \\
Lasso randomly selects one among collinear features & Group constraint: group correlated features, encourage joint selection \\
Inconsistent signs within group causes model instability & Voting determines dominant sign, penalizes inconsistent coefficients \\
\hline
\end{tabular}
\end{center}

Both constraints operate at the \textbf{regularization penalty level}, encouraging or penalizing specific coefficient patterns by adjusting soft thresholds, without changing the basic model structure.

\end{document}
